\emph{Les parties A et B sont indépendantes.}

\subsection*{Partie A}

Le tableau suivant donne la répartition de 100 personnes ayant obtenu le bac STMG en 2020 par sexe et par mention.

\begin{center}
	\begin{tblr}{hlines, vlines, colspec={l*{3}{X[c,35mm]}}}
		& Mention B ou TB & Sans mention ou AB & Total \\
		Fille & 13 & $x$ & 52 \\
		Garçon & 7 & 41 & 48 \\
		Total & 20 & 80 & 100 \\
	\end{tblr}

	{\footnotesize\emph{Source : MENJS-Depp}}
\end{center}

\begin{enumerate}
	\item Décrire par une phrase ce que représente le nombre $x$ et déterminer sa valeur.
	\item On choisit au hasard une personne ayant obtenu le bac STMG en 2020.
	\begin{enumerate}
		\item Quelle est la probabilité qu'il s'agisse d'un garçon ?
		\item Quelle est la probabilité qu'il s'agisse d'une fille et qu'elle ait obtenu son bac avec \og mention B ou TB \fg{} ?
	\end{enumerate}
	\item On choisit au hasard une fille ayant obtenu le bac STMG en 2020.

	Est-il vrai que la probabilité qu'elle ait obtenu une \og mention B ou TB \fg{} est égale à $0,25$ ?

	Justifier votre réponse.
\end{enumerate}

\subsection*{Partie B}

À la rentrée 2020, la part des filles en classe de terminale STI2D était de 10\,\%.

On choisit au hasard et successivement deux personnes inscrites en terminale STI2D à la rentrée 2020.

On admet que :
\begin{itemize}
	\item les tirages sont indépendants ;
	\item la probabilité qu'une personne choisie au hasard soit une fille est égale à $0,1$.
\end{itemize}

Quelle est la probabilité que l'on choisisse deux garçons ? On pourra s'aider d'un arbre de probabilités.
